\chapter{Vector Spaces}
The purpose of including material from Friedberg is that Hubbard's treatment of linear algebra focuses too much on $\bb{R}^n$ and many results can be stated and proved in greater generality with less redundancy and little extra effort.
\section{Definition and Examples}
For our purposes, the field we usually use is $\bb{R}$. But $\bb{C}$ and $\bb{Q}$ are common fields also.
\begin{defn}
A \emph{vector space} $V$ over a field $F$ is a space where vector addition and scalar multiplication are defined and the result is still in $V$: for any $\mb{v}, \mb{w}\in V$ and $a,b \in F$, we have $\mb{v}+ \mb{w}, a \mb{v}\in V$. Additionally, 8 axioms must be satisfied: vector addition is commutative and associative, there is an additive identity denoted $\mb{0}$, there is an additive inverse, scalar and field multiplication are compatible $(ab) \mb{v}=a(b \mb{v})$, identity of scalar multiplication $1 \mb{v}= \mb{v}$, distributivitiy of scalar multiplication with respect to vector addition $a( \mb{v}+ \mb{w})=a \mb{v}+ a \mb{w}$, and the distributivity of scalar multiplication with respect to field addition $(a+b) \mb{v}=a \mb{v}+b \mb{v}$.\footnote{if it's not already clear from the language, elements of $V$ are generally called 'vectors', while those of $F$ are generally called 'scalars'}
\end{defn}
As the language of the definition suggests, elements of a vector space are called 'vectors' (even if they might not have both magnitudes and directions as they do in $\bb{R}^n$), while those of its field are called 'scalars.' Most of the time, the field won't be explicitly stated.
A complete description of a vector space includes not only the set $V$, but also what addition and scalar multiplication are.
The motivation for the vector space axioms is that $\bb{R}^2$ with the usual notions of addition and scalar multiplication satisfy them.
Due to the associativity of vector addition, we can unambiguously add any finite number of vectors without the use of parentheses.
\begin{example}
The set of \emph{n-tuples} $F^n$ and matrices $F^{m \times n}$ with addition and scalar multiplication occurring entry by entry is a vector space. The \emph{diagonal entries} of $A \in F^{m \times n}$ are $a_{ii}$.\\
If $S$ is a nonempty set and $F$ is a field, let $\cl{F}(S,F)$ denote the set of functions $f:S \to F$. This is a vector space with the addition and scalar multiplication of functions defined the usual way: $(f+g)(s)=f(s)+g(s)$ and $(cf)(s)=c \cdot f(s)$.\\
The set $P(F)$ of all polynomials (of any degree), with addition and scalar multiplication defined as for functions, is a vector space.\\
The set of sequences with terms in $F$ with addition and scalar multiplication defined term by term is a vector space.
\end{example}
As might already be evident, addition can't be defined arbitrarily. Eg, in $\bb{R}^2$ if $(a_1,a_2)+(b_1,b_2)$ is defined as $(a_1+b_1,a_2-b_2)$ or $(a_1+b_1,0)$, then we don't have a vector space.\footnote{in the former, commutativity fails to hold. In the latter, there fails to be an additive identity}
\begin{exer}[Elementary properties of vector spaces]
Let $V$ be a vector space, $\mb{x},\mb{y},\mb{z}\in V$, $a \in F$.
\begin{itemize}
\item Cancellation law: if $\mb{x}+ \mb{z}= \mb{y}+ \mb{z}$, then $\mb{x}= \mb{y}$.
\item The additive identity is unique.
\item $\mb{x}$ has a unique additive inverse, customarily denoted $- \mb{x}$.
\item $0 \cdot \mb{x}= \mb{0}$.
\item $(-a) \mb{x}=-(a \mb{x})=a(- \mb{x})$ is the additive inverse of $a \mb{x}$.
\item $a \mb{0}= \mb{0}$.
\end{itemize}
\end{exer}

\section{Subspaces}
In the study of any algebraic structure, it's of interest to examine subsets that possess the same structure as the set under consideration.
\begin{defn}
If $V$ is a vector space over a field $F$, then $W \subseteq V$ is a \emph{subspace} of $V$ if it's a vector space over the field $F$ with addition and scalar multiplication as defined on $V$.
\end{defn}
$V$ and $\{ \mb{0}\}$ are subspaces of $V$.
\begin{lem}\label{l:check-subspace}
$W$ is a subspace of $V$ iff $\mb{0}\in W$, and for $\mb{x}, \mb{y}\in W$ and $c \in F$ we have $\mb{x}+ \mb{y},c \mb{x}\in W$.
\end{lem}
\begin{proof}
Rating 2/1.
We need for $W$ to be closed under addition and scalar multiplication. Since $W \subseteq V$ and $V$ is a vector space, all the axioms except the existence of an additive identity and additive inverses automatically carry over to $W$. The additive identity is in fact $\mb{0}$, and the additive inverse of any $\mb{x}\in W$ is $(-1) \mb{x}$.
\end{proof}
The requirement that $\mb{0}\in W$ can be replaced with $W$ being nonempty, but to show the latter it's usually easiest to show $\mb{0}\in W$.
\begin{example}
The set $\cl{S}^n(F)$ of symmetric matrices is a subspace of $F^{n \times n}$, as are the set of diagonal matrices, or the set of matrices with vanishing \emph{trace} $\tr A:= \sum_{i=1}^na_{ii}$. The set of polynomials $P_n(F)$ of degree at most $n$ is a subspace of $P(F)$. The set $C( \bb{R})$ of continuous functions $f: \bb{R}\to \bb{R}$ is a subspace of $\cl{F}( \bb{R}, \bb{R})$.
\end{example}
\begin{lem}
Any intersection of subspaces of a vector space $V$ is a subspace of $V$. If $V_1,V_2$ are 2 subspaces of $V$, then $V_1 \cup V_2$ is a subspace iff 1 contains the other.
\end{lem}
\begin{proof}
Easy, short: use Lemma \ref{l:check-subspace}.
\end{proof}

\section{Linear combinations}
As a motivating example, given 3 noncollinear vectors $\bvec{a}, \bvec{b}, \bvec{c}\in \bb{R}^3$ (or really $\bb{R}^n$), the plane determined by them is $\{ \mb{a}+s( \mb{b}- \mb{a})+t( \mb{c}- \mb{a}): s,t \in \bb{R}\}$.
\begin{defn}
If $V$ is a vector space and $S \subseteq V$ is nonempty, then $\mb{v}\in V$ is a \emph{linear combination} of vectors in $S$ if there exists $\mb{u}_1, \ldots, \mb{u}_n \in S$ and \emph{coefficients} $a_1, \ldots,a_n \in F$ such that $\mb{v}= \sum_{i=1}^na_i \mb{v}_i$.
The \emph{span} of $S$ is the set of all linear combinations of vectors in $S$ and is denoted $\Span S$. It's convenient to define $\Span \emptyset:= \{ \mb{0}\}$. We say that $S$ (or its vectors) spans $V$ if $\Span S=V$.\footnote{Friedberg also says $S$ generates $V$ if so}
\end{defn}
Note that $\mb{u}_1, \ldots, \mb{u}_n$ need not all be distinct.
\begin{lem}\label{l:span-minimal}
If $V$ is a vector space and $W$ is a subspace of $V$, $\mb{u}_1, \ldots, \mb{u}_n \in W$ and $a_1, \ldots,a_n \in F$, then $\sum_{i=1}^na_i \mb{u}_i\in W$. If $S \subseteq V$, then $S \subseteq \Span S$, and $\Span S$ is a subspace of $V$. If $S \subseteq W$, then $\Span S \subseteq W$.
\end{lem}
\begin{proof}
$\sum_{i=1}^na_i \mb{u}_i\in W$ is basically a restatement of 'closed under addition and scalar multiplication.' It can be seen with induction.\textcolor{red}{Friedberg has a link to prove this}
The rest of the statement clearly holds when $S= \emptyset$, so assume it's nonempty. Lemma \ref{l:check-subspace} verifies that $\Span S$ is a subspace of $V$. $\Span S$ contains $S$ since for any $\mb{v}\in S$, $1 \cdot \mb{v}\in \Span S$. Finally, $\Span S \subseteq W$ by the first part ($\sum_{i=1}^na_i \mb{u}_i\in W$).
\end{proof}

\section{Linear (in)dependence}
Motivating remarks: if $V$ is a vector space over an infinite field and $W$ is a subspace of $V$, then $W$ is an infinite set unless $W= \{ \mb{0}\}$. We'd like a finite and as small as possible set $S$ to span $W$: this costs us the fewest possible computations to represent an arbitrary vector in $W$.
Given vectors $\mb{u}_1, \ldots, \mb{u}_n \in W$, to see if any vector is a linear combination of the rest, 1 approach (assuming $W$ is in Euclidean space) is to try to solve for $\mb{u}_j$ using the remaining vectors, but this could require solving $n$ linear systems.
A more economical approach is to note that some vector is a linear combination of the rest iff a nontrivial linear combination of $\mb{u}_1, \ldots, \mb{u}_n$ results in $\mb{0}$ (this is formally stated in Lemma \ref{l:LID-no-lincomb3}).
\begin{defn}
Let $V$ be a vector space, then $S \subseteq V$ (or the vectors of $S$) is \emph{linearly dependent} (LD) if there exists $\mb{u}_1, \ldots, \mb{u}_n \in S$, $a_1, \ldots,a_n \in F$, not all 0, such that $\sum_{j=1}^na_j \mb{u}_j= \mb{0}$, in which case we call $\sum_{j=1}^na_j \mb{u}_j$ a \emph{nontrivial representation} of $\mb{0}$. If $S$ (or its vectors) isn't linearly dependent, then we call it \emph{linearly independent} (LID).
\end{defn}
It's immediate that, if $\mb{0}\in S$, then $S$ is linearly dependent, $\emptyset$ is vacuously linearly independent. Note that in the above definition, we can without loss of generality allow all of $a_1, \ldots,a_n$ to be nonzero.
\begin{lem}\label{l:subset-LID}
If $V$ is a vector space and $\mb{v}\in V$ with $\mb{v}\neq \mb{0}$, then $\{ \mb{v}\}$ is linearly independent.\\
If $S_1 \subseteq S_2 \subseteq V$, then the linear dependence of $S_1$ implies the linear dependence of $S_2$. The contrapositive is that the linear independence of $S_2$ implies the linear independence of $S_1$.
\end{lem}
\begin{proof}
Easy, short.\\
For the 1st part, if $a \mb{v}= \mb{0}$ for some $a \neq 0$, then $\mb{v}=a^{-1}a \mb{v}=a^{-1}\mb{0}= \mb{0}$.\\
For the 2nd part, a nontrivial representation of $\mb{0}$ using a finite set of vectors from $S_1$ would also be a nontrivial representation of $\mb{0}$ using a finite set of vectors from $S_2$.
\end{proof}
The next result helps us get at the idea of a 'minimal' spanning set.
\begin{lem}\label{l:span-LD}
If $V$ is a vector space, $S \subseteq V$ is linearly independent, and $\mb{v}\in V$, then $S \cup \{ \mb{v}\}$ is linearly dependent iff $\mb{v}\in \Span S$. 
\end{lem}
\begin{proof}
Easy, short.\\
"$\leftarrow$" If $\mb{v}\in \Span S$, then we can pick $\mb{u}_1, \ldots, \mb{u}_n \in S$, and $a_1, \ldots,a_n \in F$ such that $\mb{v}= \sum_{j=1}^na_j \mb{u}_j$. Then $\mb{v}- \sum_{j=1}^na_j \mb{u}_j$ is a nontrivial representation of $\mb{0}$ using a finite set of vectors in $S \cup \{ \mb{v}\}$.
"$\to$" If $S \cup \{ \mb{v}\}$ is linearly dependent, pick $\mb{u}_1, \ldots, \mb{u}_n \in S$ and $a,a_1, \ldots,a_n \in F$, all nonzero, such that $a \mb{v}+ \sum_{j=1}^na_j \mb{u}_j= \mb{0}$. Note that the $\mb{v}$ must be included since the linear independence of $S$ precludes a nontrivial representation of $\mb{0}$ using only vectors in $S$. Then $\mb{v}=a^{-1}\sum_{j=1}^na_j \mb{u}_j \in \Span S$.
\end{proof}
\begin{lem}\label{l:LD-lincomb}%Not yet in Friedberg
$\mb{v}_1, \ldots, \mb{v}_k$ in a vector space $V$ are linearly independent iff none of the $\mb{v}_i$ is a linear combination of the others iff none of them is a linear combination of the preceding.\\
\end{lem}
\begin{proof}
Easy, short. The proof is easier if we negate the 3 pieces.\\
"$2 \to 3$" is obvious. "$3 \to 1$" If the vectors are linearly dependent, e can find coefficients, not all 0, such that $\sum_{i=1}^ka_i \mb{v}_i= \mb{0}$. Let $m$ be the largest index where $a_m \neq 0$. Then clearly $\mb{v}_m$ is a linear combination of the preceding.
"$1 \to 2$" If $\mb{v}_1= \sum_{i=2}^ka_i \mb{v}_i$, then $\mb{v}_1- \sum_{i=2}^ka_i \mb{v}_i= \mb{0}$, with not all the coefficients 0 (eg, the 1 in front of $\mb{v}_1$).
\end{proof}

\section{Basis and dimension}
\begin{defn}
A \emph{basis} for a vector space $V$ is a linearly independent subset of $V$ that spans $V$.
\end{defn}
$\emptyset$ is vacuously a basis for $\{ \mb{0}\}$.
\begin{lem}\label{l:basis-unique-lincomb}
If $V$ is a vector space, then $\beta:= \{\mb{v}_1, \ldots, \mb{v}_n \}\subseteq V$ are linearly independent iff each $\mb{w}\in V$ can be written as a linear combination of them in at most 1 way. $\beta$ is a basis for $V$ iff each $\mb{w}\in V$ can be uniquely expressed as a linear combination of them.
\end{lem}
\begin{proof}
Easy, short.\\
Linearly independent part "$\leftarrow$" if $\mb{0}= \sum_{j=1}^n0 \mb{v}_j$ is the only way to write $\mb{0}$ as a linear combination of vectors in $\beta$, then it's linearly independent. "$\to$" If $\sum_{j=1}^na_j \mb{v}_j= \sum_{j=1}^nb_j \mb{v}_j$, then $\sum_{j=1}^n(a_j-b_j) \mb{v}_j= \mb{0}_n$. Linear independence now gives us that each $a_j=b_j$.\\
The basis part just combines the linear independence part with the fact that $\beta$ spans $V$ iff each $\mb{w}\in V$ may be written as a linear combination of vectors in $\beta$. 
\end{proof}
This is perhaps the most useful property of bases, and makes them the building blocks of vector spaces. Since any $\mb{w}\in V$ may be uniquely identified by the coefficients $a_1, \ldots,a_n$ if $\mb{w}= \sum_{j=1}^na_j \mb{v}_j$, this suggests that $V$ is 'like' $F^n$: we show this formally later.
\begin{lem}\label{l:reduce-span-to-basis}
If a vector space $V$ is spanned by a finite set $S$, then some subset of $S$ is a basis for $V$.
\end{lem}
\begin{proof}
Rating 2/2.\\
If $S= \emptyset$ or $\{ \mb{0}\}$, then $\beta= \emptyset$ works as a basis for $V= \{ \mb{0}\}$. Else, $S= \{ \mb{v}_1, \ldots, \mb{v}_n \}$ consists of nonzero vectors in $V$. Starting with $\beta= \emptyset$, we sequentially add $\mb{v}_j$ to $\beta$ if doing so maintains the linear independence of $\beta$. Thus, $\beta$ is linearly independent by construction.
Since $S$ is finite, this process must eventually stop, and it does so in 1 of 2 ways. The 1st possibility is $\beta=S$, in which case it's both linearly independent and spans $V$, so $\beta$ is the desired basis.
The 2nd possibility is that we can't add another $\mb{v}_j$ to $\beta$ without breaking linear independence. We need only show $S \subseteq \Span \beta$, for then $V= \Span S \subseteq \Span \beta$, by Lemma \ref{l:span-minimal}. For $j=\ot{n}$, either $\mb{v}_j \in \beta$, or if not, then by construction, $\beta \cup \{ \mb{v}_j \}$ is linearly dependent for any $\mb{v}_j$ not already in $\beta$, so $\mb{v}_j \in \Span \beta$ by Lemma \ref{l:span-LD}. This shows $S \subseteq \Span \beta$.
\end{proof}
The proof also suggests a way to reduce a spanning set into a basis. Later, we will see an easier way to do so by solving a linear system.
Friedberg is primarily concerned with vector spaces that have finite bases. This result identifies a large class of vector spaces where this is the case.
\begin{prop}[Replacement Theorem]\label{p:replacement}
If $V$ is a vector space spanned by a set $G$ of size $n$, and $L \subseteq V$ is a linearly independent set of size $m$, then $m \leq n$, and there exists a set $H \subseteq G$ of size $n-m$ such that $L \cup H$ spans $V$.
\end{prop}
\begin{proof}
Rating 3/2.\\
The idea is to do induction on $m$, though it might seem counterintuitive (if $m-1 \leq n$, how do we know $m \leq n$?). The base case is $m=0$, in which case $L= \emptyset$, and we can take $H=G$.
If $L= \{ \mb{v}_1, \ldots, \mb{v}_m \}$, then $L':= \{ \mb{v}_1, \ldots, \mb{v}_{m-1}\}$ is still linearly independent by Lemma \ref{l:subset-LID}, so the inductive hypothesis says that $m-1 \leq n$, and there exists $H'= \{ u_1, \ldots,u_{n-m+1}\}\subseteq G$ such that $L' \cup H'$ spans $V$.
In particular, $\mb{v}_m \in \Span(H' \cup L')$, so we have $\mb{v}_m= \sum_{j=1}^{m-1}a_j \mb{v}_j+ \sum_{j=1}^{n-m+1}b_j \mb{u}_j$. Note that both $n-m+1>0$ and some $b_j \neq 0$, lest $\mb{v}_m \in \Span L'$, which would then imply $L$ is linearly dependent by Lemma \ref{l:span-LD}.
Without loss of generality, say $b_1 \neq 0$. Put $H:= \{ \mb{u}_2, \ldots, \mb{u}_{n-m+1}\}$. Then $\mb{u}_1 \in \Span(L \cup H)$, as are all the vectors in $L' \cup H$. Thus $V= \Span(L' \cup H') \subseteq \Span(L \cup H)$ by Lemma \ref{l:span-minimal}.
\end{proof}
Friedberg calls the corollaries of this Proposition among the most important results of the chapter.
\begin{corollary}\label{c:bases-n-same}
If $V$ is a vector space that has a basis $\beta$ of $n$ vectors, then all linearly independent subsets of $V$ have at most $n$ vectors, and all bases for $V$ have $n$ vectors.
\end{corollary}
\begin{proof}
Rating 2/1.\\
Let $\gamma$ be a linearly independent subset of $V$. We don't a priori know that $\gamma$ is finite. Assume for contradiction that $\gamma$ has more than $n$ vectors. Then we can pick $S \subseteq \gamma$ having $n+1$ vectors.
Then $S$ is linearly independent by Lemma \ref{l:subset-LID}. Since $\beta$ spans $V$, Proposition \ref{p:replacement} then says $n+1<n$, which is a contradiction.
Thus $\gamma$ has at most $n$ vectors. If we further know that $\gamma$ is a basis, then switching the roles of $\gamma$ and $\beta$, we also see that assuming $\gamma$ has fewer than $n$ vectors results in a contradiction, so it in fact has exactly $n$ vectors.
\end{proof}
This result asserts that the number of vectors in any basis for a vector space is an intrinsic property:
\begin{defn}
A vector space $V$ is \emph{finite-dimensional} if it has a basis consisting of $n$ vectors, for some $n \in \bb{Z}^+$, which we call the \emph{dimension} of $V$, and write $\dim V=n$. A vector space that isn't finite-dimensional is \emph{infinite-dimensional}, write $\dim V= \infty$.
\end{defn}
\begin{example}\label{ex:dim-basis}
We have $\dim \{ \mb{0}\}=0$.
If $\bvec{e}_j \in F^n$ has jth entry 1 and all other entries 0, then $\bvec{e}_1, \ldots, \bvec{e}_n$ is the \emph{standard basis} for $F^n$, hence $\dim F_n=n$.
If $E_{ij}\in F^{m \times n}$ has ij entry 1 and all other entries 0, then $\{ E_{ij}:i=\ot{m},j=\ot{n}\}$ is a basis for $F^{m \times n}$, hence $\dim F^{m \times n}=mn$.\\
$\{ E_{11}, \ldots, E_{nn}\}$ is a basis for the diagonal matrices in $F^{n \times n}$, which thus has dimension $n$. $\{ E_{ij}+E_{ji}:i= \ot{n},j=i, \ldots,n \}$ is a basis for the symmetric matrices in $F^{n \times n}$, which has dimension $n(n+1)/2$.
The dimension of a vector space does, however, depend on the field of scalars. For instance, the vector space $V= \bb{C}$ over the field $F= \bb{C}$ has $\dim \bb{C}=1$ (a basis is $\{ 1\}$), but the same vector space over the field $F= \bb{R}$ has $\dim \bb{C}=2$ (a basis is $\{ 1,i\}$).
\end{example}
\begin{corollary}\label{c:extend-LID2basis}
Let $V$ be a vector space with $\dim V=n \in \bb{Z}^+$.
\begin{itemize}
\item If $S$ is a finite set that spans $V$, then $S$ has at least $n$ vectors. If $S$ has exactly $n$ vectors, then it's a basis.
\item If $L \subseteq V$ is a linearly independent, then $L$ has at most $n$ vectors, and there is a basis $\beta$ for $V$ such that $L \subseteq \beta$. In particular, if $L$ has exactly $n$ vectors, then $L$ is a basis for $V$.
\end{itemize}
\end{corollary}
\begin{proof}
Easy, short.\\
By Lemma \ref{l:reduce-span-to-basis}, we know that some $H \subseteq S$ is a basis for $V$, and by Corollary \ref{c:bases-n-same}, we know that $H$ has $n$ vectors. Thus, $S$ has at least $n$ vectors, and if it had exactly $n$ vectors, then $H=S$.\\
Let $L$ have $m$ vectors. If $\beta'$ is a basis for $V$ (which exists by definition of a finite-dimensional vector space), then it has $n$ vectors, and $m \leq n$ by Corollary \ref{c:bases-n-same}.
Further, we can pick $H \subseteq \beta'$ with $n-m$ vectors such that $L \cup H$ spans $V$, by Proposition \ref{p:replacement}. Since $L \cup H$ has $n$ vectors, we know it's a basis by part 1. If $m=n$, then $H= \emptyset$, so $L$ itself spans $V$, hence is a basis for $V$.
\end{proof}
\begin{lem}\label{l:polynomial-std-basis}
If $p_j(x):=x^j$ for $j \in \bb{Z}_{\geq 0}$, then $\{ p_0, \ldots, p_n \}$ is the \emph{standard basis} for $P_n(F)$, so $\dim P_n(F)=n+1$, and $P(F)$ is infinite-dimensional, with basis $\{ p_j:j \in \bb{Z}_{\geq 0}\}$.
\end{lem}
\begin{proof}
Easy, short.\\
To show $\{ p_0, \ldots,p_n \}$ is linearly independent and spans $P_n(F)$ is an exercise in Friedberg. If $\sum_{j=0}^na_jx^j=0$ for all $x \in F$, then each $a_j=0$ since Definition \ref{d:polynomial} says polynomials are equal if their coefficients match.
It's also clear that any polynomial in $P_n(F)$ (resp $P(F)$) can be written as a linear combination of $\{ p_0, \ldots,p_n \}$ (resp $\{ p_j:j \in \bb{Z}_{\geq 0}\}$).
If $S \subseteq \{ p_j:j \in \bb{Z}_{\geq 0}\}$ is finite, then there exists $N \in \bb{Z}^+$ such that $S \subseteq \{ p_j:j=0, \ldots,N \}$, the latter of which we just showed is linearly independent, so $S$ itself must be by Lemma \ref{l:subset-LID}. 
We can see that there are arbitrarily large independent subsets of $P(F)$, eg $\{ p_j:j=0, \ldots,n \}$ for any $n \in \bb{Z}^+$. We thus have that $P(F)$ is infinite dimensional by Corollary \ref{c:bases-n-same} (if $\dim P(F)=n \in \bb{Z}^+$, then any linearly independent subset must have at most $n$ vectors, precluding the existence of, say, a linearly indepedent subset of size $n+1$).
\end{proof}
At this point we're unable to prove that every infinite-dimensional vector space has a basis (note that, by definition, a finite-dimensional vector space has a  basis).
\begin{lem}\label{l:subspace-dim}
If $V$ is a vector space with $\dim V=n \in \bb{Z}^+$, and $W \subseteq V$ is a subspace, then $\dim W \leq \dim V$: in particular, $W$ is finite-dimensional. Any basis for $W$ can be extended into one for $V$. If $\dim W= \dim V$, then $W=V$.
\end{lem}
\begin{proof}
Rating 2/1: we can't a priori assume $W$ is finite-dimensional, 'obvious' as it is.\\
If $W= \{ \mb{0}\}$, then its basis is $\emptyset$, so it's finite-dimensional, and $\dim W=0 \leq n$.
Else, there exists $\mb{x}_1 \in W$ (any $\mb{x}_1 \neq \mb{0}$ will do) such that $\{ \mb{x}_1 \}$ is linearly independent . At step $k \in \bb{Z}^+$, $\{ \mb{x}_1, \ldots, \mb{x}_{k-1}\}$ is linearly independent. If there exists $x_k \in W$ such that $\{ \mb{x}_1, \ldots, \mb{x}_k \}$ is linearly independent, move onto step $k+1$.
This process must stop at step $k \leq n$ since we can't have a linearly independent subset of $W \subseteq V$ with more than $n$ vectors by Corollary \ref{c:bases-n-same}.
If once we have $\beta:= \{ \mb{x}_1, \ldots, \mb{x}_k \}$, we'd have that for all $\mb{w}\in W$, $\beta \cup \{ \mb{w}\}$ is linearly dependent, then $\mb{w}\in \Span \beta$ by Lemma \ref{l:span-LD}. This shows that $\beta$ is a basis for $W$. Thus, $\dim W=k \leq n$.
If $\beta'$ is any basis for $W$, then it's a set of linearly independent vectors in $V$, so it can be extended to a basis of $V$ by Corollary \ref{c:extend-LID2basis}.
If $k=n$, then $\beta$ is a set of $n$ linearly independent vectors in $V$, so it must be a basis for $V$ by Corollary \ref{c:extend-LID2basis}, ie $V \subseteq \Span \beta =W$ by Lemma \ref{l:span-minimal}.
\end{proof}
\begin{lem}
Let $V$ be a vector space with $\dim V=n \in \bb{Z}^+$, $S$ a set that spans $V$. Then $S$ has at least $n$ vectors, and there exists $\beta \subseteq S$ such that $\beta$ is a basis for $V$.
\end{lem}
\begin{proof}
Easy, short: the proof closely mimics that of Lemma \ref{l:subspace-dim}.\\
If $V= \{ \mb{0}\}$, then $S$ vacuously has at least 0 vectors, and $\beta:= \emptyset \subseteq S$ is a basis for $V$. Else, there exists $\mb{x}_1 \in S$ (any $\mb{x}_1 \neq \mb{0}$ will do) such that $\{ \mb{x}_1 \}$ is linearly independent.
At step $k \in \bb{Z}^+$, $\{ \mb{x}_1, \ldots, \mb{x}_{k-1}\}$ is linearly independent. If there exists $x_k \in S$ such that $\{ \mb{x}_1, \ldots, \mb{x}_k \}$ is linearly independent, move onto step $k+1$.
This process must stop at step $k \leq n$ since we can't have a linearly independent subset of $S \subseteq V$ with more than $n$ vectors by Corollary \ref{c:bases-n-same}.
Once the process stops, we claim $\beta:= \{ \mb{x}_1, \ldots, \mb{x}_k \}$ is a basis for $V$. Indeed, it's linearly independent by construction, and for all $\mb{w}\in S$, we have $\beta \cup \{ \mb{w}\}$ is linearly dependent, so $\mb{w}\in \Span \beta$ by Lemma \ref{l:span-LD}. This shows $S \subseteq \Span \beta$, hence $V= \Span S \subseteq \Span \beta$ by Lemma \ref{l:span-minimal}.
Thus, we know $k=n$ by Corollary \ref{c:bases-n-same}, so $\beta \subseteq S$ implies $S$ has at least $n$ (possibly infinite) vectors.
\end{proof}
This result is a stronger version Lemma \ref{l:reduce-span-to-basis}, where we have removed the restriction that $S$ must be finite to begin with.\\
To give a geometric interpretation, in $\bb{R}^2$, the only subspace $W$ with $\dim W=0$ is $\{ \bvec{0}_2\}$, ie the origin. If $\dim W=1$, then $W= \{ a \bvec{w}:a \in \bb{R}\}$ for some nonzero $\bvec{w}\in \bb{R}^2$: $W$ is a line thru the origin. If $\dim W=2$, then $W= \bb{R}^2$: the whole Euclidean plane.
If $W$ is a subspace of $\bb{R}^3$, then $\dim W=0$ implies $W$ is the origin and $\dim W=1$ implies a line thru the origin as before. $\dim W=2$ implies a plane thru the origin, and $\dim W=3$ implies $W= \bb{R}^3$.
\begin{prop}\label{p:Lagrange-poly}
Let $F$ be an infinite field with $c_0, \ldots,c_n \in F$ distinct. For $j=0, \ldots,n$, put $S(j):= \{0, \ldots,n \}\setminus \{ j\}$, $f_j(x):= \prod_{k \in S(j)}(x-c_k)/(c_j-c_k)$, which are \emph{Lagrange polynomials} associated with $c_0, \ldots,c_n$.
Then $\{ f_0, \ldots,f_n \}$ form a basis for $P_n(F)$, and any $g \in P_n(F)$ can be written $g= \sum_{j=0}^ng(c_j)f_j$, the \emph{Lagrange interpolation formula}. Conversely, for all $d_0, \ldots,d_n \in F$, there exists unique $g \in P_n(F)$ such that $g(c_j)=d_j$ for $j=0, \ldots,n$. In particular, if $g \in P_n(F)$ and each $g(c_j)=0$, then $g \equiv 0$.
\end{prop}
\begin{proof}
Easy, short.\\
Note $f_j(c_k)=1(j=k)$. If $\sum_{j=0}^na_jf_j \equiv 0$, then $0= \sum_{j=0}^na_jf_j(c_k)=a_k$ for $k=0, \ldots,n$, so $\{ f_0, \ldots,f_n \}$ are $n+1$ linearly independent vectors. Since we also have $\dim P_n(F)=n+1$ by Lemma \ref{l:polynomial-std-basis}, we know $\{ f_0, \ldots,f_n \}$ is in fact a basis for $P_n(F)$ by Corollary \ref{c:extend-LID2basis}.
Hence, any $g \in P_n(F)$ may be written $g= \sum_{j=0}^nb_jf_j$. Note $g(c_k)= \sum_{j=0}^nb_jf_j(c_k)=b_k$ for $k=0, \ldots,n$, so we in fact have $g= \sum_{j=0}^ng(c_j)f_j$.
\end{proof}
The process of estimating intermediate values of a variable from known values is \emph{interpolation}. For any $k>10$, say, Lagrange interpolation is a dreadful way to recover the polynomial: Bernstein polynomials, developed in the 20th century, are far better.
As an example of how they fail, consider interpolating $x \mapsto 1/(x^2+0.1)$ in the interval $[-1,1]$, sampled at 21 points that are evenly spaced by 0.1. Near the endpoints, Lagrange interpolation does a terrible job.
Interpolation is fundamental in the communications and computer industry, where a function that fits a dataset, or (in the other direction) a way to encode a function as $k$ sampled values is desired.

\section{Maximal Linearly Independent Subsets*}
The goal here is to try, to the extent possible, to extend results from the previous section to infinite-dimensional vector spaces, chiefly to prove that every vector space has a basis.
This is important for in the case of infinite dimensions since it's often hard to construct an explicit basis. As an example, $\bb{R}$ is a vector space over the field $\bb{Q}$, but there's no obvious way to construct a basis for it.
Many proofs in the previous section relied on induction, which will no longer be adequate in the case of infinite dimension.
\begin{defn}
Let $\cl{F}$ be a family of sets. $M \in \cl{F}$ is \emph{maximal} (with respect to set inclusion) if $M' \in \cl{F}$ and $M \subseteq M'$ implies $M=M'$, ie $M$ is contained in no set of $\cl{F}$ except itself.\\
The family of sets $\cl{C}$ is a \emph{chain} or \emph{tower} if $A,B \in \cl{C}$ implies $A \subseteq B$ or $B \subseteq A$.
\end{defn}
\begin{example}
For a nonempty set $S$, define its \emph{power set} $\cl{P}(S)$ to be the family of all subsets of $S$. Then $S$ is maximal in $\cl{P}(S)$. If $S$ and $T$ are disjoint and nonempty, then both are maximal in $\cl{P}(S) \cup \cl{P}(T)$.\\
If $S$ is an infinite set and $\cl{F}$ is the family of finite subsets of $S$, then $S$ has no maximal element: indeed, for any $A \in \cl{F}$, there exists $s \in S \setminus A$, so $A \subseteq A \cup \{ s \} \in \cl{F}$, hence $A$ can't be maximal.
For $n \in \bb{Z}^+$, put $A_n:= \{\ot{n}\}$. Then $\{ A_n:n \in \bb{Z}^+ \}$ is a chain.
\end{example}
\begin{axiom}[Hausdorff Maximal Principle]
Let $\cl{F}$ be a family of sets. If for all chains $\cl{C}\subseteq \cl{F}$ there exists $B \in \cl{F}$ such that $A \subseteq B$ for all $A \in \cl{C}$, then $\cl{F}$ has a maximal element.
\end{axiom}
The Hausdorff Maximal Principle is logically equivalent to the \emph{Axiom of Choice}, and it motivates reformulating the definition of a basis in terms of a maximal property.
\begin{defn}
Let $V$ a vector space, $S \subseteq V$. A \emph{maximal linearly independent subset} (MLID) $B$ of $S$ is linearly independent, and $A \subseteq S$ is linearly independent and $B \subseteq A$ implies $A=B$.
\end{defn}
\begin{example}
Any subset of $\{ \mb{e}_1, \mb{e}_2, \mb{e}_1+\mb{e}_2\} \subseteq F^n$ of size 2 is maximally linearly independent, so maximal linearly independent subsets need not be unique.
\end{example}
\begin{prop}
Let $V$ be a vector space and $S$ span $V$. Then $\beta$ is a maximal linearly independent subset of $S$ iff it's a basis for $V$.
\end{prop}
\begin{proof}
Easy, short.\\
"$\leftarrow$" $\beta$ is linearly independent, and for all $\mb{v}\in S \subseteq V$, we have that $\beta \cup \{ \mb{v}\}$ is linearly dependent by Lemma \ref{l:span-LD} (since $\beta$ spans $V$), so $\beta$ can't be a proper subset of any linearly independent $A \subseteq S$.
"$\to$" $\beta$ is linearly independent, for contradiction assume $S \nsubseteq \Span \beta$. Then for any $\mb{v}\in S \setminus \Span \beta$, we have that $\beta \cup \{ \mb{v}\} \subseteq S$ is linearly independent, contradicting that $\beta$ is maximally linearly independent.
Since $S \subseteq \Span \beta$, it follows from Lemma \ref{l:span-minimal} that $V= \Span S \subseteq \Span \beta$.
\end{proof}
\begin{thm}\label{t:basis-exists}
If $V$ is a vector space and $L \subseteq V$ is linearly independent, then there exists a maximally linearly independent $\beta \subseteq V$ such that $L \subseteq \beta$. In particular, every vector space has a basis.
\end{thm}
\begin{proof}
Rating 1/2.\\
Put $\cl{F}:= \{ A \subseteq V:A \T{ is linearly independent and }L \subseteq A \}$. To apply the Hausdorff Maximal Principle, we need to show that if $\cl{C}\subseteq \cl{F}$ is a chain, then there exists $U \in \cl{F}$ such that $B \subseteq U$ for all $B \in \cl{C}$.
If $\cl{C}= \emptyset$, then any $U \in \cl{F}$ works, eg $U=L$, else put $U:= \cup_{B \in \cl{C}}B$. We claim $U \in \cl{F}$: clearly $L \subseteq U$. To show it's linearly independent, fix $\mb{u}_1, \ldots,\mb{u}_n \in U$ and suppose $\sum_{j=0}^na_j \mb{u}_j= \mb{0}$.
Each $\mb{u}_j \in B_j$ for some $B_j \in \cl{C}$. Since $\cl{C}$ is a chain, we have that $B':= \cup_{j=1}^nB_j \in \cl{C}$. Thus, $B'$ is linearly independent, and $\mb{u}_1, \ldots, \mb{u}_n \in B'$, hence each $a_j=0$. The Hausdorff Maximal Principal now applies to give the desired conclusion.
\end{proof}
\begin{exer}
\begin{itemize}
Let $V$ be a vector space.
\item The vector space $\bb{R}$ over the field $\bb{Q}$ is infinite-dimensional, which can be deduced from the existence of a transcendental number, eg $\pi$.
\item (Generalize Lemma \ref{l:subspace-dim}) If $W$ is a subspace of $V$, then any basis of $W$ is a subset of some basis of $V$.
\item (Generalize Lemma \ref{l:basis-unique-lincomb}, Friedberg has online solution) $\beta$ is a basis for $V$ iff for all $\mb{v}\in V$, there exists unique $\mb{u}_1, \ldots, \mb{u}_n \in \beta$, $c_1, \ldots,c_n \in F$, all nonzero, such that $\mb{v}= \sum_{j=1}^nc_j \mb{u}_j$.
\item (Generalize Lemma \ref{l:reduce-span-to-basis}) If $S_1 \subseteq S_2 \subseteq V$, $S_1$ is linearly independent and $S_2$ spans $V$, then there exists a basis $\beta$ for $V$ such that $S_1 \subseteq \beta \subseteq S_2$.\footnote{Hint: mimic the proof of Theorem \ref{t:basis-exists} to $\cl{F}$, which consists of linearly independent $L$ such that $S_1 \subseteq L \subseteq S_2$}
\item (Generalize Proposition \ref{p:replacement}) If $\beta$ is a basis for $V$ and $L \subseteq V$ is linearly independent, then there exists $H \subseteq \beta$ such that $L \cup H$ is a basis for $V$.
\end{itemize}
\end{exer}
%()9